% ============================================================
\chapter{Étude et réalisation du Sprint 4}
\label{chap:sprint4}
% ============================================================

\section{Introduction}
\label{sec:intro_sprint4}

Ce chapitre présente l'étude et la réalisation du
\textbf{quatrième sprint} de la plateforme FinBank. Après la mise en
place des opérations bancaires, ce sprint enrichit la plateforme par
un assistant bancaire conversationnel et une base de connaissances
interrogeable.

Les travaux réalisés concernent la détection d'intention, la gestion
des sessions et des messages, la recherche augmentée par récupération
\textbf{(RAG)}, ainsi que la gestion des documents utilisés comme
sources de connaissance. Une fonctionnalité d'extraction de données
à partir d'un texte ou d'un fichier est également intégrée.

\section{Backlog du Sprint}
\label{sec:backlog_sprint4}

Le tableau~\ref{tab:backlog_sprint4} présente les principales tâches
réalisées durant le Sprint 4.

\begin{table}[htbp]
\centering
\caption{Backlog du Sprint 4}
\label{tab:backlog_sprint4}
\renewcommand{\arraystretch}{1.3}
\begin{tabular}{|c|p{10cm}|c|}
\hline
\textbf{ID} & \textbf{Tâches} & \textbf{Estimation} \\
\hline
S4-T1 & Mettre en place la détection d'intention des demandes utilisateur & 3 jours \\
\hline
S4-T2 & Développer la création et le suivi des sessions conversationnelles & 3 jours \\
\hline
S4-T3 & Développer l'envoi, la conservation et la consultation des messages & 3 jours \\
\hline
S4-T4 & Intégrer le service de dialogue de l'assistant bancaire & 4 jours \\
\hline
S4-T5 & Mettre en place la recherche dans la base de connaissances RAG & 4 jours \\
\hline
S4-T6 & Gérer les documents de connaissance et leur réindexation & 3 jours \\
\hline
S4-T7 & Ajouter l'extraction d'informations depuis un texte ou un fichier & 2 jours \\
\hline
S4-T8 & Intégrer et vérifier l'interface de conversation & 2 jours \\
\hline
\multicolumn{2}{|r|}{\textbf{Total}} & \textbf{24 jours} \\
\hline
\end{tabular}
\end{table}

\section{Conception}
\label{sec:conception_sprint4}

La conception du Sprint 4 décrit les interactions entre l'utilisateur,
l'assistant bancaire et les services de traitement du langage. Elle
prend également en compte le rôle de l'administrateur, qui gère les
documents constituant la base de connaissances.

\subsection{Diagramme de cas d'utilisation}
\label{subsec:usecase_sprint4}

L'utilisateur authentifié peut créer une session, envoyer des
messages, consulter la conversation et poser une question à
l'assistant. Le système détecte l'intention et recherche, lorsque cela
est nécessaire, des informations dans la base de connaissances RAG.
L'administrateur peut gérer, activer, désactiver et réindexer les
documents de cette base.

\begin{figure}[htbp]
\centering
\IfFileExists{Figures/usecase_sprint4.png}{%
  \includegraphics[width=\textwidth,height=0.78\textheight,keepaspectratio]{Figures/usecase_sprint4.png}%
}{%
  \fbox{\parbox[c][5cm][c]{0.88\textwidth}{\centering
  Espace réservé au diagramme de cas d'utilisation du Sprint 4.}}
}
\caption{Diagramme de cas d'utilisation du Sprint 4}
\label{fig:usecase_sprint4}
\end{figure}

\subsection{Diagramme de classes du Sprint 4}
\label{subsec:classes_sprint4}

Le diagramme de classes présente les entités dédiées à l'assistant :
\texttt{Session}, \texttt{Message}, \texttt{BankingAgent} et
\texttt{RagDocument}. Une session appartient à un utilisateur et peut
être associée à un agent bancaire. Les messages sont rattachés à la
session et conservent notamment l'expéditeur, l'intention détectée et
les informations relatives à l'utilisation du RAG.

\begin{figure}[htbp]
\centering
\IfFileExists{Figures/classes_sprint4.png}{%
  \includegraphics[width=\textwidth,height=0.82\textheight,keepaspectratio]{Figures/classes_sprint4.png}%
}{%
  \fbox{\parbox[c][5cm][c]{0.88\textwidth}{\centering
  Espace réservé au diagramme de classes du Sprint 4.}}
}
\caption{Diagramme de classes du Sprint 4}
\label{fig:classes_sprint4}
\end{figure}

\subsection{Description textuelle : Interagir avec l'assistant}
\label{subsec:desc_assistant_sprint4}

\begin{table}[htbp]
\centering
\caption{Cas d'utilisation : Interagir avec l'assistant bancaire}
\label{tab:uc_assistant_sprint4}
\renewcommand{\arraystretch}{1.3}
\begin{tabular}{|p{3.2cm}|p{10.5cm}|}
\hline
\textbf{Élément} & \textbf{Description} \\
\hline
Titre & Interagir avec l'assistant bancaire \\
\hline
Acteur & Utilisateur (\texttt{USER}) ou Administrateur (\texttt{ADMIN}) \\
\hline
But & Obtenir une réponse ou une orientation à partir d'une demande exprimée en langage naturel. \\
\hline
Précondition & L'utilisateur est authentifié et possède une session ouverte. \\
\hline
Postcondition & Le message et la réponse de l'assistant sont conservés dans la conversation. \\
\hline
Scénario principal &
1. L'utilisateur ouvre ou crée une session. \newline
2. Il saisit une question. \newline
3. Le système enregistre le message. \newline
4. Le système détecte l'intention. \newline
5. Le système recherche un contexte pertinent. \newline
6. L'assistant produit une réponse. \newline
7. Le système affiche et conserve la réponse. \\
\hline
Scénario alternatif & Si le service intelligent est indisponible, le système retourne une erreur explicite et conserve l'état de la demande. \\
\hline
\end{tabular}
\end{table}

\subsection{Description textuelle : Gérer les documents de connaissance}
\label{subsec:desc_documents_sprint4}

\begin{table}[htbp]
\centering
\caption{Cas d'utilisation : Gérer les documents de connaissance}
\label{tab:uc_documents_sprint4}
\renewcommand{\arraystretch}{1.3}
\begin{tabular}{|p{3.2cm}|p{10.5cm}|}
\hline
\textbf{Élément} & \textbf{Description} \\
\hline
Titre & Gérer les documents de connaissance \\
\hline
Acteur & Administrateur (\texttt{ADMIN}) \\
\hline
But & Maintenir les contenus utilisés par la recherche RAG. \\
\hline
Précondition & L'administrateur est authentifié et autorisé. \\
\hline
Postcondition & Le document est créé, modifié, supprimé ou réindexé selon l'opération demandée. \\
\hline
Scénario principal &
1. L'administrateur ouvre la gestion de la base de connaissances. \newline
2. Il saisit ou modifie le titre et le contenu. \newline
3. Le système valide les données. \newline
4. Le système enregistre le document. \newline
5. Le document peut être activé ou réindexé. \\
\hline
Scénario alternatif & Si l'acteur n'est pas administrateur ou si le contenu est invalide, l'opération est refusée. \\
\hline
\end{tabular}
\end{table}

\subsection{Diagramme de classes participantes : Conversation}
\label{subsec:participants_chat_sprint4}

Ce diagramme représente les classes intervenant dans le traitement
d'un message : l'interface de conversation, le contrôleur, le service
de chat, la détection d'intention, la recherche RAG et les entités
persistantes.

\begin{figure}[htbp]
\centering
\IfFileExists{Figures/participants_chat_sprint4.png}{%
  \includegraphics[width=\textwidth,height=0.78\textheight,keepaspectratio]{Figures/participants_chat_sprint4.png}%
}{%
  \fbox{\parbox[c][5cm][c]{0.88\textwidth}{\centering
  Espace réservé au diagramme de classes participantes de la conversation.}}
}
\caption{Diagramme de classes participantes -- Conversation}
\label{fig:participants_chat_sprint4}
\end{figure}

\subsection{Diagramme de séquence : Question avec recherche RAG}
\label{subsec:seq_rag_sprint4}

Le diagramme de séquence décrit l'envoi d'une question, la détection
de son intention, la recherche des passages pertinents et la
construction du contexte transmis à l'assistant.

\begin{figure}[htbp]
\centering
\IfFileExists{Figures/seq_rag_sprint4.png}{%
  \includegraphics[width=\textwidth,height=0.82\textheight,keepaspectratio]{Figures/seq_rag_sprint4.png}%
}{%
  \fbox{\parbox[c][5cm][c]{0.88\textwidth}{\centering
  Espace réservé au diagramme de séquence d'une question RAG.}}
}
\caption{Diagramme de séquence -- Question et recherche RAG}
\label{fig:seq_rag_sprint4}
\end{figure}

\section{Réalisation}
\label{sec:realisation_sprint4}

Cette section présente les composants développés pour fournir une
expérience conversationnelle persistante et exploiter les documents de
connaissance.

\subsection{Détection d'intention}
\label{subsec:intention_sprint4}

L'endpoint \texttt{/ai/intent} reçoit le texte de la demande et
retourne une intention ainsi que les informations associées lorsque
la détection est disponible. Le résultat peut être utilisé par le
service de conversation pour orienter le traitement de la demande.
La solution peut s'appuyer sur le service de détection configuré dans
le backend, avec une solution de repli lorsque le fournisseur
d'intelligence artificielle n'est pas disponible.

\begin{figure}[htbp]
\centering
\IfFileExists{Figures/swagger_ai_sprint4.png}{%
  \includegraphics[width=0.9\textwidth]{Figures/swagger_ai_sprint4.png}%
}{%
  \fbox{\parbox[c][4cm][c]{0.85\textwidth}{\centering
  Espace réservé à la capture Swagger des endpoints d'intelligence artificielle.}}
}
\caption{Endpoints de détection d'intention et de recherche RAG}
\label{fig:swagger_ai_sprint4}
\end{figure}

\subsection{Sessions et messages}
\label{subsec:sessions_messages_sprint4}

Le \texttt{SessionController} permet de créer, consulter, ouvrir et
fermer les sessions. Il permet également de vérifier leur validité et
de fermer toutes les sessions d'un utilisateur. Le
\texttt{MessageController} assure l'enregistrement des messages, la
consultation paginée et l'affichage de la conversation dans l'ordre
chronologique.

Chaque message conserve son contenu, son type d'expéditeur, son
intention éventuelle, ses paramètres et les informations liées à
l'utilisation du RAG. Cette traçabilité facilite la restitution de la
conversation et l'analyse des échanges.

\begin{figure}[htbp]
\centering
\IfFileExists{Figures/chat_sessions_sprint4.png}{%
  \includegraphics[width=0.9\textwidth]{Figures/chat_sessions_sprint4.png}%
}{%
  \fbox{\parbox[c][4cm][c]{0.85\textwidth}{\centering
  Espace réservé à la capture de l'interface des sessions et de la conversation.}}
}
\caption{Gestion des sessions et des messages}
\label{fig:chat_sessions_sprint4}
\end{figure}

\subsection{Recherche augmentée par récupération (RAG)}
\label{subsec:rag_sprint4}

L'endpoint \texttt{/ai/ask} reçoit une question, vérifie qu'elle n'est
pas vide, puis recherche un nombre limité de résultats pertinents dans
la base de connaissances. Les résultats sont assemblés dans un
contexte et retournés avec la question, l'intention détectée, le
niveau de confiance et l'indication de disponibilité des résultats.

Cette approche permet de fournir des réponses fondées sur les
documents enregistrés dans FinBank plutôt que sur une réponse isolée
du modèle conversationnel.

\begin{figure}[htbp]
\centering
\IfFileExists{Figures/rag_search_sprint4.png}{%
  \includegraphics[width=0.9\textwidth]{Figures/rag_search_sprint4.png}%
}{%
  \fbox{\parbox[c][4cm][c]{0.85\textwidth}{\centering
  Espace réservé à la capture de la recherche RAG et de ses sources.}}
}
\caption{Recherche dans la base de connaissances RAG}
\label{fig:rag_search_sprint4}
\end{figure}

\subsection{Gestion et réindexation des documents}
\label{subsec:documents_rag_sprint4}

Les documents de connaissance sont gérés par les endpoints
\texttt{/admin/rag-documents}. L'administrateur peut les lister, les
créer, les modifier, les supprimer et demander leur réindexation.
Chaque opération vérifie explicitement le rôle administrateur. Le
contenu, la source et l'état actif du document sont conservés afin de
contrôler les éléments proposés lors de la recherche.

\begin{figure}[htbp]
\centering
\IfFileExists{Figures/rag_documents_sprint4.png}{%
  \includegraphics[width=0.9\textwidth]{Figures/rag_documents_sprint4.png}%
}{%
  \fbox{\parbox[c][4cm][c]{0.85\textwidth}{\centering
  Espace réservé à la capture de l'administration des documents RAG.}}
}
\caption{Administration des documents de connaissance}
\label{fig:rag_documents_sprint4}
\end{figure}

\subsection{Extraction d'informations depuis un document}
\label{subsec:extraction_document_sprint4}

L'endpoint \texttt{/ai/extract-document} accepte un texte ou un fichier
multipart. Le service d'extraction identifie les champs disponibles et
retourne le type de document ainsi qu'un indicateur précisant si des
données ont été trouvées. Cette fonctionnalité complète l'assistant en
facilitant la lecture de contenus transmis par l'utilisateur.

\section{Conclusion}
\label{sec:conclusion_sprint4}

Ce chapitre a présenté l'étude et la réalisation du Sprint 4 de
FinBank. La conception a été décrite à travers les diagrammes de cas
d'utilisation, de classes, de classes participantes et de séquence.

La réalisation a permis d'intégrer la détection d'intention, les
sessions et les messages conversationnels, la recherche RAG ainsi que
la gestion des documents de connaissance. Ces composants constituent
la base de l'assistant bancaire qui sera exploité par les interfaces et
les workflows présentés dans les chapitres suivants.

